
% =========================================================
% BÀI 2: Thu thập và tải dữ liệu
% =========================================================

\section{Thu thập và tải dữ liệu}

Dữ liệu được thu thập từ bộ dữ liệu thứ cấp (định dạng CSV) lưu trữ trên Google Drive. Tệp dữ liệu chứa hơn \(54\,000\) quan sát (dòng), hoàn toàn thỏa mãn điều kiện tối thiểu \(500\) quan sát của đề bài. 

Quá trình tải và trích xuất dữ liệu được tự động hóa bằng ngôn ngữ Python (sử dụng thư viện \texttt{gdown} và \texttt{pandas}). Đoạn mã cốt lõi thực hiện tải dữ liệu và khai báo biến được minh họa như sau:

\begin{verbatim}
import gdown
import pandas as pd

# Tải file từ Google Drive
file_id = '1RgrGtIgDJQiuYjxZ3WSxESXdaZztoj9T'
download_path = 'dataset.csv'
gdown.download(id=file_id, output=download_path, quiet=False)

# Đọc dữ liệu
df = pd.read_csv(download_path)
print(f"Kích thước dữ liệu: {df.shape[0]} dòng, {df.shape[1]} cột.")

# Khai báo biến
X_feature = ['gateway_rtt_avg_ms']
target_y = 'rtt_avg_ms'
\end{verbatim}

Sau khi tải thành công, nhóm tiến hành kiểm tra sự tồn tại của các cột \texttt{gateway\_rtt\_avg\_ms} và \texttt{rtt\_avg\_ms} để đảm bảo dữ liệu sẵn sàng cho các bước tiền xử lý và phân tích khám phá (EDA) tiếp theo.